\documentclass[12pt]{article}
\usepackage[margin = 1in]{geometry} % set margins
\usepackage{amsmath, % math formatting
            amsfonts, % math fonts
            amssymb,  % extends math symbols
            mathtools, % additional math formatting
            bbm,
            amsthm
}
\usepackage{comment, % comments out large sections
            setspace, % line spacing
            hyperref % hyperlinks
}
% set spacing
\onehalfspacing
\setcounter{secnumdepth}{0}
% theorem environments (global numbering)
\newtheorem{theorem}{Theorem}
\newtheorem{lemma}[theorem]{Lemma}
\newtheorem{prop}[theorem]{Proposition}
\newtheorem{corollary}[theorem]{Corollary}
\newtheorem{remark}[theorem]{Remark}
\newtheorem{definition}[theorem]{Definition}
% =======================================
%             MACROS
% =======================================
% Math operators
\DeclareMathOperator{\E}{\mathbb{E}}
\newcommand{\R}{\mathbb{R}}
% Distribution / functions
\newcommand{\PhiStd}{\Phi}          % standard normal CDF
\newcommand{\phiStd}{\phi}          % standard normal PDF
\newcommand{\ex}[1]{\exp\!\left(#1\right)}  % exp with nice spacing
\newcommand{\logg}[1]{\log\!\left(#1\right)}% log with nice spacing
% Vectors / profiles
\newcommand{\vprof}{\mathbf v}
\newcommand{\aprof}{\mathbf a}
\newcommand{\vminus}{\mathbf v_{-i}}
\newcommand{\aminus}{\mathbf a_{-i}}
% Winner
\newcommand{\winner}[1]{w\!\left(#1\right)} % w(a)
\newcommand{\GE}{\text{GE}}
% Reduced-form objects
\newcommand{\Sagg}{S} % use \Sagg for the symbol
% Opponent strength aggregator S(v_-i, a_-i)
\newcommand{\Sof}[2]{\Sagg\!\left(#1,#2\right)} % S(v_-i,a_-i)
% Primary nomination probability p_a(v;S)
% Usage: \pnom{a}{v}{S}
\newcommand{\pnom}[3]{p_{#1}\!\left(#2;#3\right)}
% Closed-form nomination probability for generic a in {0,1}
% Usage: \pnomrf{a}{v}{S}
\newcommand{\pnomrf}[3]{\frac{\ex{#2+#1}}{\ex{#2+#1}+#3}}
% Shorthands p1 and p0 as functions of (v;S)
\newcommand{\pOne}[2]{p_{1}\!\left(#1;#2\right)}
\newcommand{\pZero}[2]{p_{0}\!\left(#1;#2\right)}
% Explicit closed forms:
\newcommand{\pOneRF}[2]{\frac{\ex{#1+1}}{\ex{#1+1}+#2}}
\newcommand{\pZeroRF}[2]{\frac{\ex{#1}}{\ex{#1}+#2}}
% Primary proportional benefit G(v;S) = p1/p0
\newcommand{\Gfun}[2]{G\!\left(#1;#2\right)}
% General proportional cost C(v;delta) = Phi(v)/Phi(v-delta)
\newcommand{\Cfun}[2]{C\!\left(#1;#2\right)}
% Delta incentive object: Delta(v;S) = p1*Phi(v-delta) - p0*Phi(v)
\newcommand{\Deltainc}[2]{\Delta\!\left(#1;#2\right)}
% General-election win probability
% Usage: \Pge{v}{delta}{a} -> Phi(v - delta a)
\newcommand{\Pge}[3]{\PhiStd\!\left(#1-#2\,#3\right)}
% Useful comparisons / relations
\newcommand{\iffc}{\;\Longleftrightarrow\;}

% ---------------------------------------------
\title{Crowded Primaries and Weakened Nominees: Behavioral Distortion in a Two-Stage Electoral Contest}
\author{Katherine Parslow}
\date{Last Updated: \today}

\begin{document}
\maketitle

\section*{Model \& Notation}

\subsection*{Setup}

$N \ge 2$ candidates compete for a party's nomination.
Each candidate $i$ has valence $v_i \in \R$, drawn i.i.d.\ from distribution $F$
with strictly positive density $f$ on $\R$.
The valence profile $\vprof = (v_1,\dots,v_N)$ is drawn by nature at the start of the game
and is observed by all candidates.

Candidates simultaneously choose an action $a_i \in \{0,1\}$.
Taking action ($a_i = 1$) boosts a candidate's relative standing in the primary
but harms their viability in the subsequent general election.
Let $\aprof = (a_1,\dots,a_N)$ denote the action profile.

\subsection*{Primary Election}

Candidates are ranked according to their \emph{primary index}
\begin{equation}
    q_i^P = v_i + a_i + \varepsilon_i,
    \label{eq:primary-index}
\end{equation}
where $\varepsilon_i \sim \text{EV1}$ are i.i.d.\ candidate-level shocks.
The nominee, denoted $w$, is the candidate with the highest primary index.

Define the \emph{opponent-strength aggregator} for candidate $i$ as
\[
    \Sof{\vminus}{\aminus} = \sum_{j \neq i} \ex{v_j + a_j}.
\]
Integrating over the EV1 shocks, the probability that candidate $i$ wins the nomination is
\begin{equation}
    \Pr(w = i \mid \vprof, \aprof)
        = \pnomrf{a_i}{v_i}{\Sof{\vminus}{\aminus}}.
    \label{eq:primary-prob}
\end{equation}
We write this more compactly as $\pnom{a_i}{v_i}{\Sof{\vminus}{\aminus}}$,
with the special cases $\pOne{v_i}{\Sof{\vminus}{\aminus}}$ when $a_i = 1$
and $\pZero{v_i}{\Sof{\vminus}{\aminus}}$ when $a_i = 0$.

\subsection*{General Election}

The nominee $w$ advances to the general election.
Their \emph{general-election index} is
\begin{equation}
    q_w^{\GE} = v_w - \delta\, a_w + \eta_w,
    \label{eq:ge-index}
\end{equation}
where $\eta_w \sim N(0,1)$ is an idiosyncratic shock independent of the primary shocks,
and $\delta > 0$ captures the cost of primary electioneering.
Normalizing the out-party opponent's index to $0$, the probability that nominee $w$ wins the general election is
\begin{equation}
    \Pr(\GE\text{ win} \mid v_w, a_w) = \Pge{v_w}{\delta}{a_w}.
    \label{eq:ge-prob}
\end{equation}

\subsection*{Payoffs}

Candidates are office-motivated and receive a payoff of $1$ from holding office and $0$ otherwise.
Because holding office requires winning both elections, candidate $i$'s expected utility is
\begin{equation}
    u_i(\vprof, \aprof)
        = \pnom{a_i}{v_i}{\Sof{\vminus}{\aminus}} \cdot \Pge{v_i}{\delta}{a_i}.
    \label{eq:utility}
\end{equation}

\subsection*{Timing}

The timing of the game is as follows:
\begin{enumerate}
    \item Nature draws the valence profile $\vprof = (v_1,\dots,v_N)$;
    \item Candidates observe $\vprof$ and simultaneously choose $a_i \in \{0,1\}$ for all $i$;
    \item The primary election is held and the party nominee $w$ is selected according to \eqref{eq:primary-prob};
    \item The general election is held and payoffs are determined according to \eqref{eq:ge-prob}.
\end{enumerate}

\subsection*{Notation Summary}

\begin{center}
\begin{tabular}{cl}
    \hline
    \textbf{Symbol} & \textbf{Description} \\
    \hline
    $N$ & Number of candidates \\
    $v_i$ & Valence of candidate $i$; drawn i.i.d.\ from $F$ \\
    $\vprof = (v_1,\dots,v_N)$ & Valence profile \\
    $a_i \in \{0,1\}$ & Candidate $i$'s primary-electioneering action \\
    $\aprof = (a_1,\dots,a_N)$ & Action profile \\
    $\delta > 0$ & Cost of taking the primary action in the general election \\
    $\Sof{\vminus}{\aminus}$ & Opponent-strength aggregator $\sum_{j\neq i}\exp(v_j+a_j)$ \\
    $\pnom{a}{v}{S}$ & Nomination probability for action $a$, valence $v$, aggregator $S$ \\
    $\Phi(\cdot)$ & Standard normal CDF; used for GE win probability \\
    \hline
\end{tabular}
\end{center}

\end{document}